\section{Introduction}\label{sec:introduction}

A mixed-integer approximation of a nonlinear graph can spend its resources
in different ways. It can introduce many short interval selectors, use a
few integer variables with large ranges, add continuous lifting variables,
or encode delicate geometry in long rational coefficients. Counting only
the intervals of a particular approximation does not determine the minimum
number of integer variables available to all formulations. This paper asks
for that minimum while retaining the entire input-output graph.

For a continuous map $F:D\to\R^m$ and a prescribed error body $K$, an
admissible projected relaxation $R$ must satisfy
\[
 \{(x,F(x)):x\in D\}\subset R
 \subset\{(x,w):x\in D,\ w-F(x)\in K\}.
\]
We minimize the number of general integer coordinates in an arbitrary
convex lift of such an $R$, denoting the result by $\pconv$.
The comparison quantity $\pbin$ restricts the lift to linear constraints
and binary coordinates. Continuous dimension is finite but unrestricted
in both minima, and coefficients may be real. Thus $\pconv\le\pbin$.
A constructive count $\pout$ is reported separately when the formulation
has polynomial rational encoding and can be produced in polynomial time.
These definitions measure formulation resources; they do not assert that
the resulting mixed-integer problem is easy to solve.

The lower bounds examine exact graph contacts in the lift. Witnesses in
one integer parity class have admissible midpoints. Covariance and volume
bounds on these contacts give precision-dependent lower bounds even when
the original witness sets are not measurable or the lift is not closed.
Other examples require thirds combinations or convexity inside a fixed
binary section. The upper bounds encode carefully chosen local graph bands
and share integer indices across nonlinear expressions. This contact-and-band
approach links integer dimension to the geometry and algebra of the map.

The principal quadratic result identifies an exact algebraic precision
coefficient. If $\mathcal H$ is the real span of the Hessians of a fixed
quadratic system on a full-dimensional box, then
\begin{equation}\label{eq:intro-ncrank}
 \pconv(\eps),\ \pbin(\eps)
 =\tfrac12\ncrank(\mathcal H)\log_2(1/\eps)+O(1).
\end{equation}
Noncommutative rank can exceed the maximum rank of a scalar linear
combination. For the three outputs of a cross product on six inputs,
every scalar combination has Hessian rank at most four, but the joint
precision coefficient is three. Consequently the largest scalar rank is
not the invariant for a quadratic system. The proof combines symmetric
shrunk-subspace geometry and operator capacity with contact covariance
bounds and a common binary discretization. Scalar quadratic rank,
one-sided inertia, and fractional vertex-cover laws for bilinear
interaction graphs appear as related exact cases.

After normalizing the input box to $[0,1]^n$, the finite-tolerance
determinant program
\[
 \max\{\det P:0\preceq P\preceq I,
       \ \operatorname{tr}(H_jPH_jP)\le\eps_j^2\ \forall j\}
\]
produces a benchmark within $O(n\log(n+1))$ of both minima.
Its matrix constraints need not be Euclidean convex, but the problem has
useful geodesic convexity. We give finite certificates, a rational
polynomial-time construction, output-body and nonlinear-input reductions,
and sharper positive semidefinite block and diagonal results. Correlated
input domains require their own geometry; a thin-domain example shows
why a containing-box calculation cannot simply replace the actual domain.
We also prove a computational barrier to uniformly improving additive
count approximation below every fixed sublinear power of input dimension,
unless $P=NP$.

Beyond quadratics, the analysis separates local curvature from effective
finite-accuracy shape. Smooth maps admit local/global Hessian-rank bounds;
constant scalar Hessian rank and positive perspectives give exact rates.
For one continuous convex scalar graph, a three-piece chord refinement
puts binary linear dimension within two of $\pconv$. An accuracy-dependent
truncated curvature mass supports certified rational knot construction.
A hybrid of exact greedy segmentation and curvature integration gives a
polynomial rational MILP within eleven binaries for every dense convex
polynomial. Positive separable systems permit coefficient allocation,
feature-curve lower bounds and a double-logarithmic degree overhead, with
an implementation polynomial in sparse binary-exponent encoding. Pure
powers admit degree-independent finite comparisons and separate compact
constructions for binary rational exponents greater than one.

For several convex outputs, a basis chosen from original outputs preserves
nonnegative chord gaps. Its curvature rank replaces output count in the
additive comparison. Positive facet scalarizations and positive polar
spanners extend this result to unconditional error bodies, including
bodies supplied only by a strong separation oracle. The final formulation
uses an explicit rational inner band and makes no oracle calls during
optimization. One common basis across all input coordinates also gives
separable multivariate bounds. The roles of convexity, input dimension,
and error geometry are tested by explicit separations: fixed-degree convex
product graphs have exact counts $n$ and $\lceil n\log_2 3\rceil$;
a nonconvex scalar polynomial family has two general integers and growing
binary dimension; and a tilted two-output body transfers that gap to
componentwise convex polynomials while its condition number grows.

Table~\ref{tab:result-map} gives a route through the main statements.
The symbols for different ranks and allocation benchmarks are defined
locally; they are not interchangeable invariants.
\begin{table}[htbp]
\centering\small
\begin{tabular}{@{}>{\raggedright\arraybackslash}p{.25\textwidth}>{\raggedright\arraybackslash}p{.44\textwidth}>{\raggedright\arraybackslash}p{.23\textwidth}@{}}
\toprule
Setting & Main conclusion & Location\\
\midrule
Quadratic systems & Exact coefficient $\ncrank(\mathcal H)/2$ & \cref{thm:ncrank}\\
Finite quadratic accuracy & Covariance count within $O(n\log(n+1))$; rational construction & \cref{thm:finite-covariance,thm:rational-finite}\\
Smooth and perspective maps & Local/global bounds; exact constant-rank scalar rate & \cref{thm:smooth-ranks,thm:constant-rank,prop:perspective}\\
Convex scalar graphs & Finite two-bit and dense compiled eleven-bit comparisons & \cref{lem:scalar-chords,thm:convex-hybrid}\\
Positive polynomial systems & Allocation bounds; dense and sparse double-logarithmic degree overhead & \cref{thm:positive-loglog}\\
Convex vector and separable maps & Logarithmic curvature-rank overhead; explicit facet and oracle bands & \cref{thm:vector-rank,thm:vector-oracle-rank,thm:separable-vector-rank}\\
Count separations & Exact convex product counts; nonconvex degree gaps & \cref{thm:exact-box-gap,thm:nonconvex-gap}\\
Rational encoding & Four-bit root graph: long MILP, short MISOCP & \cref{thm:root-encoding,thm:root-conic-separation}\\
\bottomrule
\end{tabular}
\caption{Main results and the models in which they are proved. Finite count
comparisons do not alone imply short rational descriptions.}
\label{tab:result-map}
\end{table}

\subsection{Relationship to established methods}

Minimum integer dimension for convex representations and the midpoint
obstruction are established in mixed-integer convex representability
\citep{lubin2022}. Binary-code embeddings for finite polyhedral unions
\citep{vielma2018}, product envelopes \citep{mccormick1976}, and shared
binary discretizations \citep{beach2024,beach2024b,teles2013} provide
formulation tools. Our comparisons minimize over every admissible set
between the exact graph and its error tube, and use these tools for
matching or quantitatively bounded upper constructions.

Noncommutative rank, shrunk subspaces, operator capacity and their rational
algorithms are imported algebraic results \citep{ggow2020,iqs2018}.
Their connection to the precision coefficient in \eqref{eq:intro-ncrank}
requires the contact-volume and shared-coordinate proofs given here.
Geodesic optimization on positive definite matrices and convex-body
separation/optimization are likewise established
\citep{sra2013,zhang2016,gls1981}. We specify the rational approximation,
feasibility repairs and residual certificates needed for the formulation
claims, rather than assuming exact real arithmetic in a bit-complexity
statement.

Scalar midpoint/chord comparisons and adaptive convex approximation have
close precedents \citep{simchowitz2018,codsi2021}. The rational compiler
uses established continuous encodings of Boolean circuits
\citep{avis2019} and discrete-function products \citep{adams2012}.
Shared breakpoints for several same-input outputs are explicit in
\citet{lyu2026}. Barycentric spanners \citep{awerbuch2004}, approximate
optimization within spanner construction \citep{plevrakis2020}, and
positive-polar oracle equivalence \citep{gls1981} underlie the vector
transfer. The claims here are the complete graph-error and integer-count
comparisons with their stated rational interfaces; these approximation,
selection and oracle primitives are not presented as new methods.

Finally, the cap-set estimate is the theorem of \citet{ellenberg2017};
our use of it is a reduction from graph-contact residues. Bernstein
approximation and difference-of-convex decomposition are classical as
well. The explicit examples distinguish which hypotheses support a
count comparison. We make no claim that the sufficient constants are
optimal, that every open smooth or vector case is characterized, or that
small integer dimension implies an ideal linear relaxation. In particular,
a universal additive binary/general-integer gap for one-input convex
vectors with growing output dimension and box error remains open here.

The remaining sections proceed from contact geometry and quadratic laws,
through finite rational constructions and scalar shape, to coupled vector
theory and the separations. Relative error and the root-power encoding
barrier are included because they show why the accuracy model and coefficient
representation cannot be omitted from an integer-dimension statement.
